
\chapter*{Abstract}

Music information retrieval often requires integrating signal-derived, algorithmic, and symbolic information across different representations. Existing visualization tools typically keep notation separate from physical time, while composites that combine them are assembled manually. This thesis presents LayerIt, a Python-based object-oriented framework for composing independent visual representations into unified figures on a shared performance-time axis. Built on a modular layer-composition model, LayerIt treats audio plots, algorithmic outputs, symbolic scores, and other data sources as independent components that can be flexibly combined and regenerated. Figures are exported as SVG, preserving the score's MEI structure while keeping added components identifiable. The thesis asks how heterogeneous MIR visualizations can be combined in a common score-aligned representation, to what extent modular layers and scriptable composition support their generation, modification, and regeneration, and how such visualizations can support qualitative inspection of beat-tracking outputs and annotation data. LayerIt is validated qualitatively through a set of demonstrations built on real score and performance data, showing the composition of score, audio, and algorithmic representations on a common performance-time axis, while establishing that the usefulness of the resulting visualizations depends on the quality of the available score-performance alignment. More broadly, the project establishes a reusable, extensible foundation for composable score-aligned MIR visualizations.

\vspace{\fill}
\begin{flushleft}
{\Large \textbf{Keywords}:} 
Music Information Retrieval \(\bullet\)
Visualization \(\bullet\)
Audio-to-Score Alignment \(\bullet\)
Beat Tracking \(\bullet\)
SVG \(\bullet\)
Object-Oriented \(\bullet\)
Python
\end{flushleft}
\vspace*{24mm}

\acresetall %% to reset the acronym usage
